\part{Choix généraux d'organisation des développements}
\label{part:dev_gene}

Cette partie traite du cadre général des développement dans AIDA, c'est à dire
l'écriture (ou la modification) de programme ou fonction d'usage général, de
documentation, mais pas l'utilisation d'AIDA pour effectuer tel ou tel
traitement.

Les sources sont stockés dans un système de gestion
de source (\texttt{git}).

%% problème avec dev.bib  : remplacement de \cite{Miller1997} par l'url seule
Le processus de compilation/installation est basé sur l'utilisation récursive
de make\footnote{\label{make}Cette technique est controversée, notamment par
Peter Miller, qui explique pourquoi les make récursifs ne permettent pas de
voir l'arbre complet des dépendances, et obligent soit à en négliger certaines,
soit à faire du travail en double
(\url{http://aegis.sourceforge.net/auug97.pdf}).  AIDA n'est pas un gros projet
(au sens de P. Miller) et on parvient à faire une compilation ordonnée
(bibliothèques avant programmes) qui évite les dépendances complexes.}
sur toute l'arborescence AIDA.


\section{Utilisation des makefile}
\label{sec:make}

\subsection{Cibles usuelles}
Chaque makefile devra définir au moins trois cibles  :

\begin{description}
  \item[clean] : supprime tous les fichiers générés au cours du processus de
    construction.
  \item[depend] : calcule les dépendances qui peuvent être construites
    automatiquement.
  \item[install] : construit le résultat attendu (bibliothèque, programme,
    document, ...) dans le répertoire actuel puis copie le résultat vers un
    répertoire
    (\texttt{\$AIDA\_HOME/lib, \$AIDA\_HOME/bin, \$AIDA\_HOME/doc, ...}).
\end{description}
\bigbreak

Il est souhaitable que le makefile définisse la cible \texttt{\textbf{all}} :
\begin{description}
  \item[all] : construit le résultat attendu (bibliothèque, programme,
    document, ...) dans le répertoire actuel (le plus souvent, "all" est la
    cible par défaut).
\end{description}
\bigbreak

Une cinquième cible, \texttt{\textbf{check}}, qui exécute des tests, est pour
l'instant recommandée mais pas encore systématique.


\subsection{Configuration et portabilité}
Pour éviter d'effectuer au sein de chaque makefile une reconnaissance de la
machine hôte, et une adaptation à l'environnement local (versions et options
des compilateurs et autres utilitaires), ces réglages sont regroupés dans un
fichier unique \path{$AIDA_HOME/config/local_cfg.mk}, et tous les makefiles
incluent ce fichier de configuration grâce à la ligne :
\smallbreak
\begin{boxedverbatim}
-include $(AIDA_HOME)/config/local_cfg.mk 
\end{boxedverbatim}
\medbreak

Le fichier \path{$AIDA_HOME/config/local_cfg.mk} contient toutes les
définitions de macros pouvant être différentes selon l'environnement. C'est un
lien sur :\\
\path{AIDA_HOME/local/make/Linux_$DISTRIB_$DISTRIB_RELEASE_$ARCH.auto.mk}. \\

Le fichier
\path{AIDA_HOME/local/make/Linux_$DISTRIB_$DISTRIB_RELEASE_$ARCH.auto.mk} et
le lien \path{$AIDA_HOME/config/local_cfg.mk} sont générés automatiquement par
l'installateur (voir : \path{$AIDA_HOME/scr/script/set_make_config.sh})

\subsection{Quelques points techniques spécifiques}
\begin{itemize}
\item Pour compatibilité avec
un système de gestion de sources,
les dépendances générées automatiquement par
makedepend ne sont pas écrites en fin de makefile,
mais dans un fichier séparé.
%TODO : explication de la technique
\item .SECONDEXPANSION pour gmake version $\geq$ 3.81
\end{itemize}

